% Independent supplement: --main supplementary.tex --no-build-script.
\documentclass[pdflatex,sn-nature]{sn-jnl}
\usepackage{graphicx,amsmath,amssymb,amsfonts,booktabs,array,tabularx,siunitx,longtable}
\sisetup{uncertainty-mode=separate}
\renewcommand{\thesection}{S\arabic{section}}
\renewcommand{\thefigure}{S\arabic{figure}}
\renewcommand{\thetable}{S\arabic{table}}
\renewcommand{\theequation}{S\arabic{equation}}

\begin{document}
\title[Supplementary information]{Supplementary information: Manuscript title}
\author*[1]{\fnm{First} \sur{Author}}\email{first.author@example.org}
\author[1]{\fnm{Second} \sur{Author}}
\affil*[1]{\orgdiv{Department},\orgname{Institution},
  \orgaddress{\city{City},\country{Country}}}
\abstract{Replace this instructional summary with an overview of the
supplement, or remove it if the journal's current instructions require a
different opening. This document is independent of the main manuscript.}
\maketitle

\section{Supplementary methods}
Describe supporting details and derivations. All figures, tables and equations
in this document have independent S numbering. The example expression
in equation~\eqref{eq:si-residual} illustrates numbering only.
\begin{equation}
\begin{split}
 r_i&=y_i-\widehat y_i,\\
 E&=\frac{1}{n}\sum_{i=1}^{n}r_i^2.
\end{split}\label{eq:si-residual}
\end{equation}

\section{Supplementary figures}
Figure~\ref{fig:si-slot} demonstrates a separately cited supporting figure.
Use a descriptive title and explain the actual symbols and uncertainties.
\begin{figure}[htbp]
\centering
\fbox{\parbox[c][55mm][c]{0.85\linewidth}{\centering
Supplementary figure slot\\[8pt]
Replace with an original PDF or PNG.}}
\caption{\textbf{Supporting graphic.} Replace this legend with a description
of the study's actual figure. The box is a layout example and represents no
data. A supplementary figure and its legend should fit on a single page.}
\label{fig:si-slot}
\end{figure}
\clearpage

\section{Supplementary tables}
Table~\ref{tab:si-numeric} demonstrates compact numerical columns.
Table~\ref{tab:si-long} demonstrates repeated headings for a table too long
for one page. Retain only the environments needed by the study.
\begin{table}[htbp]
\centering\small
\caption{Synthetic entries used only to illustrate numerical alignment.}
\label{tab:si-numeric}
\begin{tabularx}{\linewidth}{@{}>{\raggedright\arraybackslash}X
  S[table-format=2.0] S[table-format=1.3(3)]@{}}
\toprule
Setting & {Count} & {Estimate $\pm$ SD}\\
\midrule
Example A & 24 & 0.320 +- 0.015\\
Example B & 36 & 0.415 +- 0.021\\
\bottomrule
\end{tabularx}
\par\smallskip\footnotesize\raggedright
All values are invented. Replace them with verified study results and define
the actual measure of uncertainty.
\end{table}
\clearpage

\small
\begin{longtable}{@{}p{0.15\linewidth}p{0.29\linewidth}p{0.46\linewidth}@{}}
\caption{Synthetic supporting inventory to demonstrate continuation.}
\label{tab:si-long}\\
\toprule Item & Setting & Explanation\\\midrule
\endfirsthead
\multicolumn{3}{@{}l}{Table~\thetable\ (continued)}\\[3pt]
\toprule Item & Setting & Explanation\\\midrule
\endhead
\midrule\multicolumn{3}{r@{}}{Continued on next page}\\
\endfoot
\bottomrule
\endlastfoot
S01 & Identifier & Record the stable identifier of a supporting input.\\
S02 & Version & Record the exact version used in the study.\\
S03 & Origin & Describe the provenance of the supporting material.\\
S04 & Units & Define the measurement units used by the source.\\
S05 & Inclusion & Describe the rule used to select source records.\\
S06 & Exclusion & Describe the rule used to exclude source records.\\
S07 & Partition & Explain separation of development and evaluation material.\\
S08 & Transform & Describe preprocessing without changing the original data.\\
S09 & Scaling & Define any scaling or normalisation reference.\\
S10 & Parameters & Record the numerical analysis settings.\\
S11 & Seed & Record the random seed where one is relevant.\\
S12 & Repetition & Define what counts as an independent repetition.\\
S13 & Counts & Report the exact sample count for each analysis.\\
S14 & Variability & Define the variability shown by the data.\\
S15 & Uncertainty & Define the uncertainty measure and its calculation.\\
S16 & Interval & State the interpretation of any reported interval.\\
S17 & Comparison & Identify the reference condition for a comparison.\\
S18 & Controls & Describe the controls supporting interpretation.\\
S19 & Validation & Describe independent checks supplied by the study.\\
S20 & Sensitivity & Explain a sensitivity analysis where it was performed.\\
S21 & Software & Record software versions and relevant options.\\
S22 & Environment & Describe relevant computational environment details.\\
S23 & Precision & Preserve the justified source precision in displayed values.\\
S24 & Rounding & Explain any rounding applied for presentation.\\
S25 & Aggregation & Describe how observations were combined.\\
S26 & Dependence & Identify repeated observations sharing a common source.\\
S27 & Labels & Check labels at the intended final figure dimensions.\\
S28 & Figure source & Retain plotting code and editable original graphics.\\
S29 & Table source & Link displayed values to verified analysis outputs.\\
S30 & Symbols & Define mathematical variables and dimensions.\\
S31 & Assumptions & State assumptions needed to use a derivation.\\
S32 & References & Check the supplementary bibliography independently.\\
S33 & Availability & Describe an accurate access route and restrictions.\\
S34 & Limitations & State limitations supported by the actual study design.\\
\end{longtable}
\normalsize

\section{Supplementary references}
The citation here tests this document's independent bibliography
\cite{wilkinson2016fair}; it is not evidence for the synthetic examples.
\bibliography{references}
\end{document}
